\documentclass[11pt]{article}

\usepackage[utf8]{inputenc}
\usepackage[T1]{fontenc}
\usepackage{charter}
\usepackage{amsmath,amssymb}
\usepackage{booktabs}
\usepackage{array}
\usepackage{graphicx}
\usepackage{xcolor}
\usepackage{microtype}
\usepackage{setspace}
\usepackage{lineno}
\usepackage{cite}
\usepackage[hidelinks]{hyperref}
\usepackage[a4paper,margin=2.5cm]{geometry}
\usepackage[font=small,labelfont=bf]{caption}
% Auto-generated by analysis/build_public_core_macros.py; do not edit.
% Source: results/public_core_evidence.json sha256=c41aff9fa9deeaac8a77d43e95b4029ca1e9ad3565f465c35e73e15741426caf
\newcommand{\DffRows}{12,651}
\newcommand{\DffTargets}{44}
\newcommand{\DsRows}{260,060}
\newcommand{\DsTargets}{58}
\newcommand{\DffRawPR}{1.834}
\newcommand{\DffResidualPR}{9.303}
\newcommand{\DsRawPR}{2.266}
\newcommand{\DsResidualPR}{18.206}
\newcommand{\DffRawMeanRtwo}{0.535}
\newcommand{\DffResidualMeanRtwo}{0.087}
\newcommand{\DsRawMeanRtwo}{0.431}
\newcommand{\DsResidualMeanRtwo}{0.038}
\newcommand{\DffAxisCosine}{0.970}
\newcommand{\DsAxisCosine}{0.969}
\newcommand{\DffCovRawPR}{1.915}
\newcommand{\DffCovResidualPR}{6.911}
\newcommand{\DsCovRawPR}{2.474}
\newcommand{\DsCovResidualPR}{8.235}
\newcommand{\DffNullMedianLow}{42.400}
\newcommand{\DffNullMedianHigh}{42.628}
\newcommand{\DsNullMedianLow}{56.481}
\newcommand{\DsNullMedianHigh}{56.612}
\newcommand{\DffMWConditionedNullPR}{34.875}
\newcommand{\DsMWConditionedNullPR}{47.039}
\newcommand{\DffMWConditionedExplainedShare}{0.238}
\newcommand{\DsMWConditionedExplainedShare}{0.251}
\newcommand{\DffMWConditionedMapRho}{0.924}
\newcommand{\DsMWConditionedMapRho}{0.718}
\newcommand{\DffMWConditionedSignAgreement}{0.834}
\newcommand{\DsMWConditionedSignAgreement}{0.727}
\newcommand{\DffMWQuartileRho}{0.205}
\newcommand{\DsMWQuartileRho}{0.351}
\newcommand{\DffMWMatchedControlRho}{0.992}
\newcommand{\DsMWMatchedControlRho}{0.994}
\newcommand{\DffMWContinuumRho}{0.945}
\newcommand{\DsMWContinuumRho}{0.988}
\newcommand{\DffSizeAxisRhoRange}{0.200--0.207}
\newcommand{\DsSizeAxisRhoRange}{0.322--0.361}
\newcommand{\DffOtherAxisRhoRange}{0.703--0.887}
\newcommand{\DsOtherAxisRhoRange}{0.787--0.831}
\newcommand{\DffProbeTwoHundred}{0.940}
\newcommand{\DffProbeFiveHundred}{0.973}
\newcommand{\DsProbeTwoHundred}{0.920}
\newcommand{\DsProbeFiveHundred}{0.965}
\newcommand{\DffProbeTwoHundredSign}{0.896}
\newcommand{\DffProbeFiveHundredSign}{0.931}
\newcommand{\DsProbeTwoHundredSign}{0.877}
\newcommand{\DsProbeFiveHundredSign}{0.920}
\newcommand{\DffScaffoldProbeTwoHundred}{0.925}
\newcommand{\DffScaffoldProbeFiveHundred}{0.958}
\newcommand{\DsScaffoldProbeTwoHundred}{0.914}
\newcommand{\DsScaffoldProbeFiveHundred}{0.958}
\newcommand{\DffRandomHoldoutTwoHundred}{0.935}
\newcommand{\DffRandomHoldoutFiveHundred}{0.968}
\newcommand{\DsRandomHoldoutTwoHundred}{0.917}
\newcommand{\DsRandomHoldoutFiveHundred}{0.960}
\newcommand{\DffSelectorRawRtwo}{0.900}
\newcommand{\DsSelectorRawRtwo}{0.698}
\newcommand{\DffSelectorResidualRtwo}{0.278}
\newcommand{\DsSelectorResidualRtwo}{0.219}
\newcommand{\DffSelectorDirectResidualRtwo}{0.278}
\newcommand{\DsSelectorDirectResidualRtwo}{0.219}
\newcommand{\DffSelectorResidualPCOneRtwo}{0.180}
\newcommand{\DsSelectorResidualPCOneRtwo}{0.065}
\newcommand{\DffSelectorDirectResidualDelta}{0.000015}
\newcommand{\DsSelectorDirectResidualDelta}{-0.000001}
\newcommand{\DffSelectorCommonResidualDelta}{0.036}
\newcommand{\DsSelectorCommonResidualDelta}{0.058}
\newcommand{\DffPanelTwoHundredDelta}{0.0153}
\newcommand{\DffPanelFiveHundredDelta}{0.0083}
\newcommand{\DsPanelTwoHundredDelta}{0.0113}
\newcommand{\DsPanelFiveHundredDelta}{0.0056}
\newcommand{\ExperimentalCrossRho}{0.848}
\newcommand{\ExperimentalCrossRawRho}{0.582}
\newcommand{\ExperimentalCrossQAP}{0.0001}
\newcommand{\DavisSplitRho}{0.757}
\newcommand{\PkistwoSplitRho}{0.915}
\newcommand{\DavisSameSupportLigands}{56}
\newcommand{\PkistwoSameSupportLigands}{154}
\newcommand{\DavisSameSupportRho}{0.307}
\newcommand{\PkistwoSameSupportRho}{0.326}
\newcommand{\DavisBroadSupportRho}{0.300}
\newcommand{\PkistwoBroadSupportRho}{0.329}
\newcommand{\HotspotLigands}{176}
\newcommand{\HotspotTargets}{21}
\newcommand{\HotspotSplitRho}{0.794}
\newcommand{\HotspotSplitRhoLow}{0.690}
\newcommand{\HotspotSplitRhoHigh}{0.854}
\newcommand{\DavisHotspotRawRho}{0.690}
\newcommand{\DavisHotspotResidualRho}{0.765}
\newcommand{\PkistwoHotspotRawRho}{0.526}
\newcommand{\PkistwoHotspotResidualRho}{0.751}
\newcommand{\PanelTransferPositiveCells}{14}
\newcommand{\PanelTransferTotalCells}{15}
\newcommand{\PanelTransferConservativeMean}{0.0607}
\newcommand{\PanelTransferJointAUC}{0.0616}
\newcommand{\PanelTransferJointP}{0.0069}
\newcommand{\PanelRankPositiveCells}{14}
\newcommand{\PanelRankConservativeMean}{0.0384}
\newcommand{\PanelRankP}{0.0056}
\newcommand{\PanelTargetLooPositive}{21}
\newcommand{\PanelTargetLooTotal}{21}
\newcommand{\PanelTargetLooAucLow}{0.0294}
\newcommand{\PanelTargetLooAucHigh}{0.0622}
\newcommand{\PanelTargetLooCellWinsLow}{11}
\newcommand{\PanelTargetLooCellWinsHigh}{15}
\newcommand{\HotspotClusterMedian}{0.0628}
\newcommand{\HotspotClusterLow}{0.0319}
\newcommand{\HotspotClusterHigh}{0.1063}
\newcommand{\HotspotClusterPositiveFraction}{1.000}
\newcommand{\DffMedianImputeResidualPR}{9.528}
\newcommand{\DffMissingAwareResidualPR}{9.663}
\newcommand{\DffCompleteCaseResidualPR}{13.970}
\newcommand{\DffEqualNResidualPRLow}{9.128}
\newcommand{\DffEqualNResidualPRHigh}{9.445}
\newcommand{\SharedExactLigands}{582}
\newcommand{\SharedConnectivityBlocks}{804}
\newcommand{\SharedTargetIdentities}{9}
\newcommand{\SharedReceptorStructures}{3}


\hypersetup{
  pdftitle={Two-way centering reveals distributed target-relative structure and limits score-surface compression in Vina docking matrices},
  pdfauthor={Adeliya Leleytner, Victor Safronov, Maxim Fedorov}
}

\doublespacing
\linenumbers
\renewcommand{\arraystretch}{1.15}

\title{Two-way centering reveals distributed target-relative structure and
limits score-surface compression in Vina docking matrices}

\author{
Adeliya Leleytner$^{1,*}$, Victor Safronov$^{1}$, Maxim Fedorov$^{1}$\\[4pt]
\small $^{1}$Institute for Information Transmission Problems of the Russian Academy of Sciences\\
\small (Kharkevich Institute), Bolshoy Karetny per. 19, Moscow 127051, Russia\\
\small $^{*}$Correspondence: \texttt{a.leleytner@gmail.com}
}
\date{}

\begin{document}
\maketitle

\begin{abstract}
\textbf{Background.}
Shared ligand effects can dominate dense reverse-docking matrices and obscure
how targets differ relative to one another. We asked what structure remains
after removing additive ligand and target effects, whether it can be compressed
from a small fixed target subset, and whether that distinction matters for
designing experimental kinase panels.

\textbf{Results.}
We audited two separately generated Vina-family matrices containing
\DffRows{} ligands by \DffTargets{} targets and \DsRows{} by \DsTargets{}.
Two-way centering increased participation-ratio dimension from \DffRawPR{} to
\DffResidualPR{} and from \DsRawPR{} to \DsResidualPR{}, while the residual
maps remained more correlated than three transformation-matched nulls and
varied across ligand-property domains. Global residual-map edge ordering was
recoverable from 200 random same-library ligands, although exact extreme-edge
recovery was lower. Selected-target reconstruction exposed a sharper boundary:
with eight observed targets, variance-weighted out-of-group $R^2$ was
\DffSelectorRawRtwo{} and \DsSelectorRawRtwo{} for raw scores but only
\DffSelectorResidualRtwo{} and \DsSelectorResidualRtwo{} for residual scores.
In a fixed 21-kinase comparison using DAVIS, PKIS2 and Anastassiadis
HotSpot profiles, panels selected from residual Vina maps without assay values
had lower signed co-response coverage loss than raw-map panels in
\PanelTransferPositiveCells{}/\PanelTransferTotalCells{} panel--size cells; the
joint normalized loss improvement was \PanelTransferJointAUC{}
(exploratory one-sided $p=\PanelTransferJointP{}$ in a post-hoc shared
target-label randomization on the fixed 21-target set, adjusted for the maximum
over three objectives). Row
centering also increased pairwise agreement among all three experimental maps,
whereas exact-support Vina--experiment map agreement remained moderate.

\textbf{Conclusions.}
The tested Vina matrices contain a dominant, readily compressible ligand axis
and a weaker distributed target-relative component. Separating these components
changes what a reduced target panel preserves. The result supports
transformation-matched spectral auditing and assay-value-blind panel design,
not affinity prediction, unseen-target retrieval or a general claim that
centering improves docking.

\medskip\noindent\textbf{Scientific Contribution.}
We connect spectral concentration to a concrete design consequence: raw Vina
scores can appear highly compressible because they preserve a shared ligand
axis, whereas residual target-relative variation is distributed across more
targets. A fixed-target, assay-value-blind comparison across three public kinase
panels tests this distinction while explicitly delimiting it from affinity
prediction and biological validation.
\end{abstract}

\noindent\textbf{Keywords:} AutoDock Vina; reverse docking; target-correlation
map; chemical applicability domain; effective dimension; panel design;
reproducibility.

\section{Introduction}

Reverse docking and multi-target virtual screening turn individual docking
calculations into a ligand-by-target matrix. Such matrices underpin target
fishing, off-target assessment and multiobjective ligand design
\cite{luo2017,lapillo2019,lee2020crds,krause2023}. They are also increasingly
available at a scale at which correlations among targets are impossible to
ignore: DOCKSTRING, for example, contains more than 15 million Vina scores and
was explicitly proposed for multiobjective optimization and transfer-learning
benchmarks \cite{dockstring2022}. If many target columns respond similarly to
the same compounds, treating them as independent objectives can duplicate
Vina-defined computation and obscure what a computational panel distinguishes.

The shared variation has familiar sources. Empirical docking scores respond to
ligand size, flexibility and contact opportunity, while receptors differ in
pocket geometry and score distribution \cite{pan2003,vigers2004,jacobsson2006}.
Target-wise and two-directional normalizations have therefore long been used in
reverse docking \cite{carta2007,sepresa2009,casey2009,luo2017}, and principal
components of docking profiles have been used as compound or protein
descriptors \cite{fukunishi2005,fukunishi2006,fukunishi2018}. Neither
normalization nor PCA is new. What remains unresolved is an operational question:
whether the target relationships inferred from a large docking matrix are
properties of the receptor panel, or properties of that panel \emph{on the
ligand library used to estimate them}.

This distinction matters because a correlation map is a population-dependent
object. For a score matrix $X\in\mathbb{R}^{N\times P}$, its target map contains
correlations between columns across the observed $N$ ligands. Changing the
ligand distribution can change those correlations even if the receptors and
scoring program are unchanged. Conversely, if a map can be estimated from a
small, representative sample of the intended library, it can guide an early
decision---for example, prioritizing a compact computational subset---without
requiring the complete $N\times P$ screen. Recovering a map and compressing the
underlying score surface are nevertheless different tasks. A panel can recover
the dominant ligand-wide score tendency while failing to preserve weaker
target-relative deviations. Global matrix geometry is also not a ligand-wise
target-ranking rule: subtracting one constant from every target score of a
ligand cannot change that ligand's target order.

Related matrix-design problems have been studied with experimental kinase
profiles: focused mini-panels can approximate a prespecified selectivity index,
compound-based networks can identify representative kinases, and informer-set
methods choose compounds for a new assay \cite{bembenek2018,zhang2019informer,
laufkotter2020}. Profile-QSAR and reverse-docking profiles instead use measured
or predicted activity vectors as descriptors for supervised prediction
\cite{martin2011pqsar,lee2012profiles}. Protein-subset selection in CANDO was
optimized against drug--indication recovery, whereas active-learning studies
have selected individual compound--assay experiments using compound or protein
features \cite{mangione2019cando,kangas2014active}. Those supervised objectives
are not interchangeable with ours. Here a fixed target subset and its
reconstruction model are learned only from a fully scored Vina pilot; utility is
then measured on ligands and omitted target columns that played no role in
selection or fitting. Experimental profiles remain an external boundary rather
than labels used to choose the Vina panel.

We address these questions with two separately generated Vina-family score
matrices, one spanning 44 secondary-pharmacology receptors and one spanning 58
DOCKSTRING targets \cite{nikitin2025,dockstring2022}. We first separate a shared
ligand-wide score axis from the residual target map and use participation-ratio
(PR) dimension as a compact re-expression of mean squared target correlation,
not as a new statistic. We then test map transport across molecular-property
domains and ask a deployment-level question: after fully scoring a pilot set,
how well can selected target scores reconstruct uncomputed raw and residual
target columns for chemically held-out ligands? Finally, three public kinase
profiles---DAVIS, PKIS2 and Anastassiadis HotSpot---provide a fixed-target test
of whether panels selected from raw or residual Vina maps cover signed
experimental co-response neighbourhoods \cite{davis2011,wu2025davis,drewry2017,
anastassiadis2011}. Experimental values are never used for Vina panel selection,
and every plausible evaluation-compound connectivity is removed from the Vina
calibration support.

The contribution is therefore an empirical and operational audit, rather than
a proposed normalization algorithm. Table~\ref{tab:scope} states the decisions
licensed by each object and the interpretations deliberately kept outside the
claim.

\begin{table}[htbp]
\centering
\caption{Objects in the audit and their interpretation boundaries.}
\label{tab:scope}
\small
\begin{tabular}{p{0.22\linewidth}p{0.34\linewidth}p{0.34\linewidth}}
\toprule
Object & Supported use & Not established here \\
\midrule
Correlation spectrum & Compact dependence audit on a stated matrix and
preprocessing contract & Algebraic rank, intrinsic molecular dimension or
docking quality \\
Residual target map & Computational redundancy and signed co-variation on a
stated ligand support & A portable protein network or true biophysical
specificity \\
Pilot panel and reconstruction & Approximate prediction of omitted Vina score
columns and a score-cell break-even within a declared ligand domain & Exact
recovery of target-specific residuals, extreme-score hit lists or transport to
an unspecified chemical population \\
Experimental panel coverage & Assay-value-blind coverage of signed co-response
neighbourhoods on the fixed 21-kinase panel & Affinity calibration, causal
validation, compound ranking or retrieval of unmeasured targets \\
\bottomrule
\end{tabular}
\end{table}

\section{Methods}

\subsection{Study design, estimands and analysis status}

This study was a secondary computational analysis of two redistributed dense
Vina-family score matrices (Docking-44 and DOCKSTRING-58), two redistributed
experimental kinase panels (DAVIS and PKIS2), and the checksum-verified official
Anastassiadis HotSpot workbook, which is fetched but not redistributed. No new
docking or experimental measurements were performed. The primary matrix-level estimand
was the Pearson target-correlation map on an explicitly stated finite ligand
support, before and after removal of additive ligand and target effects.  Map
agreement was Spearman correlation between corresponding strict-upper-triangle
edges.  The experimental boundary analysis used the same map estimand on a
fixed panel of 21 kinases and, for the DAVIS/PKIS2 same-support comparison, on
exact matched ligands. Panel-transfer utility was signed experimental
co-response coverage on those same fixed target labels; assay values were not
used to select Vina panels.

The analyses were not preregistered and no prospective power calculation was
performed.  Molecular weight was selected after preliminary inspection; its
quartile comparison, threshold sweep and continuum analysis are exploratory.
The descriptor-specificity, pilot-recovery, panel-compression,
chemical-group-holdout, experimental-reliability and three-panel transfer
analyses were post hoc revision analyses. Randomization probabilities and repeated-sample ranges are
therefore used to characterize a specified finite-support procedure, not as
confirmatory error control over the analysis history.  Ligand libraries, target
panels and assay systems are treated as fixed observed supports; no interval in
this study is a confidence interval for a superpopulation of targets.  The
cross-panel QAP is conditional on the 21 named target labels.

\subsection{Public score matrices, preprocessing and overlap audit}

Docking-44 was read from the frozen project export
\path{data/frozen/df_final_v4.csv.gz}, derived from Nikitin et al.
\cite{nikitin2025}.  It contains 12,651 ligand rows and 44 Vina-GPU~2.0 target
columns \cite{ding2023}.  Numeric score cells were clipped above at zero and
the 8,159 missing cells (1.466\% of 556,644 cells, distributed across 3,354
rows) were replaced by the mean of their target column for the primary
full-support analysis.  The frozen matrix already contained no positive value;
111 cells were exactly zero, so the clipping operation did not alter it.

DOCKSTRING-58 was read from the DOCKSTRING v1 table
\texttt{data/frozen/}\allowbreak\texttt{dockstring-}\allowbreak\texttt{dataset.tsv.gz}
\cite{dockstring2022}.  The source
contains 260,155 molecules and 58 AutoDock Vina target columns
\cite{vina2010}.  The primary support retained the 260,060 rows complete across
all 58 targets, excluding 95 rows containing any of 339 missing cells.  Of the
retained scores, 5,393 were positive and were clipped to zero.  Retaining the
positive values and retaining all 260,155 rows under target-mean or
target-median imputation were recorded preprocessing sensitivities in the
public producer.

The two matrices were generated by different source projects, but were not
treated as statistically independent samples.  Chemical overlap was audited at
three levels: full Standard InChIKey, its first 14 characters (the connectivity
block), and non-empty canonical Bemis--Murcko scaffold \cite{murcko1996}.  For
Docking-44, Standard InChIKeys were recomputed from the analysis SMILES.  For
DOCKSTRING, they were read from the frozen, row-aligned identity contract
\path{data/frozen/dockstring_identity_contract_2026-08-03.csv.gz}; this contract
parses the raw source SMILES with the pinned RDKit version and records the raw
Standard InChIKey and connectivity block without applying a fragment-parent or
uncharging transform.  Empty Murcko scaffolds were not counted in the
cross-matrix scaffold intersection.  Target identities and the curated
same-receptor-structure indicator were taken from the audited mapping table.
These audits produced the reported counts of shared
identities, connectivity blocks, scaffolds, nine mapped targets and three
shared receptor structures; they did not remove any row from the full spectral
analyses.

The analysis starts at these published, frozen score matrices.  It reruns all
matrix transformations and comparisons reported here, but does not rerun
receptor preparation, protonation, search-box definition, pose generation or
the upstream docking campaigns.  Receptor structures, configurations and
upstream commands for DOCKSTRING remain part of its source release.  Thus the
study audits dependence in released scores and does not reproduce poses,
absolute affinities or docking accuracy.

\subsection{Target maps, centring and spectral summaries}

For an $N\times P$ finite score matrix $X$, the raw map was the sample Pearson
correlation matrix of its $P$ target columns.  The residual surface was the
two-way-centred matrix
$\Gamma_{ij}=X_{ij}-\bar X_{i\cdot}-\bar X_{\cdot j}+\bar X_{\cdot\cdot}$, and
the residual map was the target correlation matrix of $\Gamma$.  Correlation
places target columns on a common unit-variance scale using sample standard
deviations (denominator $N-1$).  Because correlation already removes column
means, the raw-to-residual map change is operationally the effect of removing
each ligand's row mean.  Two-way centring imposes zero row sums and hence one
zero spectral mode.

The eigenspectrum was calculated from a symmetrized target correlation matrix;
for both $X$ and $\Gamma$, correlation construction standardizes every target
to unit sample variance after the stated transformation, so the primary
raw--residual PR contrast is normalization matched.
numerically negative eigenvalues were clipped to zero.  Participation-ratio
(PR) dimension was
\[
 r_{\mathrm{PR}}=\frac{(\sum_{k=1}^{P}\lambda_k)^2}
 {\sum_{k=1}^{P}\lambda_k^2}.
\]
For a $P$-target correlation matrix, $\sum_k\lambda_k=P$ and
$\sum_k\lambda_k^2=P+2\sum_{j<k}r_{jk}^2$, giving the exact identity
\[
r_{\mathrm{PR}}=\frac{P}{1+(P-1)\overline{r_{jk}^{2}}},
\]
where the bar averages unique off-diagonal target pairs.  PR is therefore a
monotone re-expression of mean squared off-diagonal target correlation, not an
algebraic rank or an independent test statistic.  Entropy effective rank, PC1
variance fraction and the number of
components reaching 90\% cumulative variance were retained as secondary
spectral summaries \cite{mazzucato2016,roy2007}.  For the shared-axis diagnostic,
$X$ was column-standardized with sample standard deviations, the leading
eigenvector was oriented to have positive loading sum, and its cosine with the
unit-normalized all-ones vector was calculated.

The covariance sensitivity replaced correlation by the sample covariance
matrix of $X$ or $\Gamma$ (denominator $N-1$) and applied the same eigenvalue and
PR formula.  It intentionally retained target-specific variances; no target
standardization was applied in this sensitivity.  Agreement between any two
maps was Spearman correlation between their corresponding $P(P-1)/2$ strict
upper-triangle entries, with average ranks for ties.  Sign agreement,
root-mean-square edge error and
deterministic top-edge overlaps were reported where specified below.

\subsection{Transformation-matched residual nulls}

Each matrix used one deterministic simple-random support of 12,000 ligands,
sampled without replacement (NumPy seeds 60 for Docking-44 and 61 for
DOCKSTRING-58).  The sorted zero-based row indices were encoded as little-endian
64-bit integers and SHA-256 hashed.  All four null families within a dataset
used that identical fingerprinted support and 500 replicates.

In the additive Gaussian null, the sampled surface was decomposed into its
grand mean, ligand mean minus the grand mean, and target mean minus the grand
mean.  Independent normal noise was generated for every cell, with each target
column's standard deviation set to that target's empirical two-way-residual
sample standard deviation.  Raw and re-centred residual PR were then
recalculated.  In the empirical residual-column null, each target column of the
observed two-way residual was independently permuted across ligands, after
which the simulated matrix was two-way centred again and its PR calculated.
This preserves each target's residual marginal at the permutation step, but not
cross-target correspondence or ligand-specific residual magnitude.  In the
row-norm null, independent standard-normal $P$-vectors were generated for each
ligand, their within-row means were removed, and each vector was rescaled to
the observed Euclidean norm of that ligand's residual row.  This preserves the
support, target count, zero row sums and every row norm exactly, but not target
marginals, target variances or alignment.

The molecular-weight-conditioned null operated on processed raw scores rather
than on residuals.  RDKit molecular weights were divided by empirical quantiles
into 1, 5, 10 or 20 bins; duplicate quantile edges were removed and exact ties
were assigned together with \texttt{searchsorted(side=left)}.  Within every bin
and target, raw scores were independently permuted across ligands.  Row means
were then removed and target correlations recomputed.  At the permutation step
this preserves the exact target-specific empirical score distribution within
each coarse size bin, including non-Gaussian shape, but destroys cross-target
ligand correspondence and does not preserve continuous within-bin trends or
residual row norms.  For 5, 10 and 20 bins, a size-matched random partition
used the same bin counts after randomly assigning bin labels to ligands.  Ten
bins were primary; the other counts tested granularity.  Besides residual PR,
the null recorded PC1 fraction, mean absolute and root-mean-square correlation,
maximum absolute edge, Spearman agreement with the observed map, sign agreement
and top-edge overlap.

Null summaries are Monte Carlo distributions under these four specified data-
generating mechanisms.  Central 90\% and 95\% simulation ranges were calculated
by empirical quantiles.  The lower-tail Monte Carlo probability was
$(1+\#\{r_{\mathrm{PR}}^{(b)}\leq r_{\mathrm{PR}}^{\mathrm{obs}}\})/(500+1)$.
The ranges and probabilities do not test docking accuracy, biological
correctness or transport beyond the fixed panels.

\subsection{Molecular-weight domains and chemical-support controls}

RDKit \texttt{Descriptors.MolWt} was calculated after parsing the frozen
analysis SMILES \cite{rdkit2025}.  Docking-44 used all 12,651 primary processed
rows.  DOCKSTRING used a score-blind simple-random sample of 15,000 complete
rows without replacement (primary seed 71); the complete-row index set is
fingerprinted in the result bundle.  Five additional fixed supports (seeds 72,
73, 74, 75 and 20260809) assessed DOCKSTRING support sensitivity.

For the quartile analysis, empirical 0.25 and 0.75 quantiles were calculated and
all boundary ties were retained: low rows satisfied
$\mathrm{MW}\leq q_{0.25}$ and high rows satisfied
$\mathrm{MW}\geq q_{0.75}$.  Raw and row-centred maps were estimated separately
in the two bands.  Two controls distinguished a domain change from ordinary
finite-support or group-disjoint variation.  Equal-size row-random halves were
non-overlapping; if a band contained an odd number of rows, one row from that
permutation was unused.  For the chemical-group control, complete groups were
assigned
within ten stable rank strata of group-median molecular weight.  Within each
stratum the group order was permuted, and groups were assigned greedily to the
currently smaller ligand-count half; the initial side for exact ties was random
and then alternated.  No group crossed halves.  The same row-random and
MW-stratified whole-group splits were repeated separately inside the low and
high bands.

Docking-44 grouping used the frozen source column
\texttt{Butina\_clusters}; the strict public-core producer did not regenerate
those labels, so their fingerprint recipe is not assumed from the label name.
DOCKSTRING grouping used canonical RDKit Bemis--Murcko scaffold SMILES.  Every
acyclic DOCKSTRING molecule was deliberately assigned its own source-row label
\texttt{ACYCLIC\_SINGLETON:<source-row>}, including acyclic molecules with the
same connectivity, so that it could not leak across group-disjoint halves.
This singleton rule is specific to the DOCKSTRING chemical-domain controls.

There were 200 repetitions per control.  The base seed was 20260812 for
Docking-44 and 20261812 for DOCKSTRING.  Within each dataset, low- and high-band
row-random controls used base seed plus 100 and 101, and low- and high-band
group-disjoint controls used base seed plus 200 and 201.  Central 95\% repeated-
split ranges describe composition sensitivity, not sampling confidence
intervals.

Threshold robustness used exact equal-size lower and upper rank tails at
fractions 0.10, 0.15, 0.20, 0.25, 0.30, 0.35 and 0.40.  Rows were sorted by
molecular weight and then source-row index, making boundary ties deterministic;
each tail contained $\lfloor fN\rfloor$ rows.  For each fraction, 100 equal-$N$
random-disjoint controls used seeds 20310810 (Docking-44) and 20311810
(DOCKSTRING).  The continuum analysis divided the same supports into ten
non-overlapping equal-$N$ rank windows; the highest-ranked $N\bmod 10$ remainder
rows were excluded.  Spearman correlation between the 45 pairwise differences
in median molecular weight and $1-$map agreement was descriptive because the
45 comparisons reuse ten maps.  Molecular weight is correlated with other
chemical properties, and none of these analyses identifies a causal molecular-
weight effect.

\subsection{Descriptor-axis specificity}

As a post hoc specificity analysis, the same low-versus-high support contrast
was repeated on seven RDKit axes computed directly from the frozen SMILES:
heavy-atom count (\texttt{GetNumHeavyAtoms}), molecular weight
(\texttt{Descriptors.MolWt}), Labute approximate surface area
(\texttt{CalcLabuteASA}), topological polar surface area
(\texttt{Descriptors.TPSA}), Crippen cLogP (\texttt{MolLogP}), rotatable-bond
count (\texttt{Descriptors.NumRotatableBonds}) and ring count
(\texttt{CalcNumRings}) \cite{rdkit2025}.  Docking-44 used all 12,651 primary
processed rows; DOCKSTRING used the same sorted 15,000-row complete-case sample
drawn without replacement with seed 71 as the molecular-weight controls.  Its
complete-row indices were SHA-256 fingerprinted.

For each descriptor separately, the empirical 25th and 75th percentiles were
calculated and boundary ties were retained with $z\leq q_{0.25}$ and
$z\geq q_{0.75}$.  Consequently, discrete descriptors could yield more than
25\% of the support in a tail; the actual low- and high-tail counts were
recorded and used for their matched controls.  Raw and row-centred target maps
were estimated on both extremes and compared by upper-triangle Spearman
correlation; sign-flip fraction and the PR of each extreme map were secondary
summaries.  All 21
pairwise Spearman correlations among the seven descriptors were also recorded
to expose their dependence.  The three labels ``size related'' (heavy-atom
count, molecular weight and Labute ASA) and ``other coarse'' were used only for
descriptive summaries.  For each distinct pair of low- and high-tail sizes, the
row-centred analysis also used 200 matched-size random-row-disjoint controls:
a seeded permutation supplied the first $N_{\mathrm{low}}$ rows to one map and
the next $N_{\mathrm{high}}$ rows to the other.  Size pairs were sorted
lexicographically and identical pairs were simulated once and reused.  The
generator seed was the dataset base plus $100{,}000$ times the zero-based size-
pair index; dataset bases were 20260821 for Docking-44 and 21260821 for
DOCKSTRING.  Control medians and empirical 2.5th and 97.5th percentiles were
descriptive finite-support comparators; no randomization probability was
assigned.  Because all seven axes were inspected post hoc and are correlated,
their contrasts are neither independent tests nor a causal decomposition of
the molecular-weight result.

\subsection{Pilot-map recovery and target-panel compression}

For each full processed source matrix, 100 simple-random pilot samples of 200
and 500 rows were drawn without replacement (seeds 20260810 for Docking-44 and
20261810 for DOCKSTRING).  The generator advanced continuously across the two
sample sizes.  Residual maps were calculated from row-centred scores.  For every
replicate, the row-disjoint complement map was computed exactly by subtracting
the pilot residual sum and cross-product from full-support sufficient
statistics; no approximation or resampling was used for the complement.

Pilot and reference maps were compared by upper-triangle Spearman correlation,
root-mean-square error and edge-sign agreement.  Strong-edge sign recovery used
all reference edges at or above the empirical 90th percentile of absolute
correlation.  Exact sets of the 10 and 25 most positive and most negative edges
were selected with stable mergesort tie handling and compared by overlap,
precision and Jaccard index.  Average-linkage clustering with optimal leaf
ordering used signed distance $1-r$ and was cut at $k=6,8,10$ for the cluster
sensitivity.

Panel compression treated either strong positive or strong negative linear
dependence as computational redundancy.  For a $k$-target panel $S$, the loss
on map $C$ was
\[
 L(S;C)=\frac{1}{P}\sum_{j=1}^{P}\min_{s\in S}\{1-|C_{js}|\}.
\]
The main decision used $k=8$; $k=6$ and $k=10$ were sensitivities.  Pilot maps
used deterministic PAM BUILD followed by SWAP, with target-name ordering to
break objective ties.  The fixed full-source reference panel was the exact
mixed-integer $k$-medoids solution with zero relative MIP gap and a negligible
index penalty used only to break exact ties.  Both panels were evaluated on the
row-disjoint complement.  The displayed difference is pilot-panel loss minus
the loss of this one fixed full-source panel; it is not regret against a panel
re-optimized on each complement.

For each dataset and each $k$, 10,000 uniformly sampled target subsets provided
the full-source random-panel comparator.  Their seeds were the dataset recovery
seed plus $10,000+k$ (20270818 and 20271818 for the main $k=8$ analyses).  To
limit repeated complement computation, the first 200 of those same fixed random
panels were evaluated on every complement.  Pilot summaries and their central
95\% ranges are empirical summaries of 100 samples from the observed library;
they neither establish a universal calibration size nor test transport to a
shifted library.

\subsection{Chemical-group-held-out pilot recovery}

To test whether within-source recovery depended on reusing chemical groups, a
separate sensitivity used five-fold \texttt{GroupKFold} with one harmonized
structure-derived grouping rule.  Each SMILES underwent RDKit cleanup,
largest-fragment selection, uncharging, isotope removal and removal of
stereochemistry.  Canonical non-isomeric Bemis--Murcko scaffold SMILES defined
cyclic groups.  Acyclic compounds were grouped by the first block of the
Standard InChIKey of the standardized parent, with canonical non-isomeric
SMILES as a declared fallback if InChI generation failed.  Docking-44 used all
12,651 rows and 5,112 groups; DOCKSTRING used the fixed 15,000-row complete-case
support drawn with seed 71 and 11,903 groups.  This Docking-44 grouping is
distinct from the frozen \texttt{Butina\_clusters} labels used in the
molecular-weight domain controls.  GroupKFold was deterministic and placed
every group wholly in one of five folds; the code asserted zero group overlap
between the four-fold calibration pool and the held-out fold.

Scores were clipped above at zero before splitting.  DOCKSTRING had no missing
score on this support.  For each Docking-44 split, target means were estimated
only from the four-fold calibration pool and applied to missing cells in both
that pool and its held-out fold.  Thus no held-out score informed imputation.
The held-out residual target map used every ligand in the fold.  From the
filled calibration pool, 25 simple-random samples without replacement were
drawn at each calibration size, 200 and 500; each row-centred map was compared
with the held-out-group map by upper-triangle Spearman correlation.

A matched random-row-holdout comparator used the same evaluation-fold size and
calibration size for every fold and repetition.  An independent permutation of
all rows supplied the random held-out rows first and the calibration pool
second; pool-only target means were again applied to both parts, and the first
$n$ pool rows formed the calibration sample.  Chemical groups were allowed to
overlap in this comparator.  Docking-44 used dataset base seed 20260805 and
DOCKSTRING used 20260806.  For one-based fold $f$ and calibration size $n$, the
chemical-group sample generator used
\texttt{dataset\_seed+10000*f+n}; the matched-random-row generator used that
seed plus 1,000,000.  Each generator advanced through the 25 replicates.

For every fold and design, achiral 2,048-bit radius-2 Morgan fingerprints were
computed on the same standardized parents.  The exact maximum Tanimoto
similarity of every held-out ligand to every ligand in its complete calibration
pool was recorded, requiring 25.6--36.0 million comparisons per fold.  The
group rule is deliberately stricter than a row split but is not a
fingerprint-cluster-disjoint split: a Tanimoto value of one denotes equality of
the declared finite bit vectors, not necessarily molecular identity, and
chemotypes with different group labels can remain similar.

The 125 values per dataset, design and calibration size were summarized by
their mean, sample standard deviation, extrema and empirical 2.5th and 97.5th
percentiles.  Paired chemical-group-minus-random-row differences used matching
dataset, fold, size and repetition indices and were summarized by their mean,
median and empirical 2.5th and 97.5th percentiles.  These ranges combine fixed-
fold composition with repeated calibration samples within one source
collection.  They are same-source holdout sensitivities, not confidence
intervals, evidence of transport to a shifted deployment library, an
equivalence test, or recovery of individual ligand scores.

\subsection{Chemical-group-held-out score-surface reconstruction}

A separate five-fold \texttt{GroupKFold} experiment tested whether a selected
target subset reconstructed held-out score surfaces. It used the same
standardized-parent chemical-group rule described above and asserted zero group
overlap. Only rows complete across all target columns entered the pilot and
evaluation sets. Within each fold, 500 pilot ligands were sampled from the four
training folds; the fifth fold remained unused for centring, scaling, target
selection, Ridge regularization or model fitting. The evaluation predictor
received only the observed raw Vina scores at the selected targets.

Panels of $k=4,6,8,10,12$ targets were selected by exact $k$-medoids on distance
$1-r^2$ from the pilot row-centred residual map. A pivoted-QR selector on the
pilot raw-score correlation representation and 50 fixed random panels per fold
and $k$ were controls. Multivariate Ridge regression predicted the omitted raw
targets; its penalty was selected from $\{0.01,0.1,1,10,100\}$ by generalized
cross-validation using pilot data only. Selected observed scores were copied
into the full predicted raw profile. That profile was then row-centred in raw
score units and standardized with pilot-only residual column scales to define
the derived residual prediction. Primary plots report variance-weighted
$R^2=1-\sum_j\mathrm{SSE}_j/\sum_j\mathrm{SST}_j$ on omitted targets; pooled
RMSE, target correlations, ligand-profile correlations, full-profile metrics,
a selected-target PC1 model and a seven-descriptor Ridge model were retained as
sensitivities.

At $k=8$, a direct-outcome sensitivity fitted the held-out residual columns
from the same selected raw target predictors, using the same pilot-only Ridge
procedure. Agreement with the derived-residual route tests whether low residual
reconstruction was mechanically introduced by predicting and then centring raw
scores. A designed-versus-random comparison was restricted to targets omitted
by both the designed panel and each of the first three fixed random panels; the
reported common-omitted contrasts are descriptive because only three random
panels were fixed for this robustness calculation.

\subsection{Experimental panels, exact ligand matching and map reliability}

The fixed target order was ABL1, AKT1, AKT2, CDK2, CSF1R, EGFR, FGFR1, IGF1R,
JAK2, KDR, KIT, LCK, MAP2K1, MAPK1, MAPK14, MAPKAPK2, MET, PLK1, PTK2, ROCK1
and SRC.  Source aliases were mapped explicitly rather than inferred from
strings.  DAVIS-Complete V3 was read from the exact CC0 TSV export
\path{data/frozen/davis_complete.tab.gz}; pivoting the selected proteins by
\texttt{drug\_name} produced a dense 72-by-21 pKd block \cite{davis2011,wu2025davis}.
PKIS2 was read from sheet \texttt{Table 4 - PKIS2 \%Inh} of the official S4
workbook \path{data/frozen/pkis2_s4.xlsx}.  Rows missing any of \texttt{Regno},
\texttt{Compound}, \texttt{Chemotype} or \texttt{Smiles} were excluded, leaving
645 rows; the 21 explicitly mapped percentage inhibition/displacement columns
were required to be numeric, dense and within $[0,100]$ \cite{drewry2017}.

Raw and row-centred target maps were computed independently within each assay
matrix. No ligand mean, target scale or other centring statistic was shared
between DAVIS, PKIS2 and HotSpot; cross-assay agreement therefore compares
separately estimated maps rather than a jointly normalized matrix.

The Anastassiadis HotSpot panel was read from Supplementary Table~3 of the
official publisher workbook \cite{anastassiadis2011}. The fetcher records the
source URL and accepts the file only at its frozen SHA-256; the workbook is not
redistributed. Percentage remaining activity was converted to
$100-\text{percentage remaining}$ without clipping. CDK2 was the per-compound
arithmetic mean of the cyclin-A and cyclin-E rows. Of 178 compounds, SB~202474
was excluded for a missing AKT1 value and VEGF Receptor~2 Kinase Inhibitor~II
for a missing MAPK14 value, leaving a dense \HotspotLigands{}-by-21 block.
Target aliases, worksheet cells, exclusions and alternative CDK2 contexts are
serialized in the result bundle.

Profiles with range no greater than $10^{-12}$ across targets were called
constant; the informative cross-panel maps therefore used 67 DAVIS and 644
PKIS2 profiles, whereas the primary split-half summaries used every released
row (72 and 645, respectively).  The two-way-centred map was primary.  Raw maps
were repeated for split-half map reliability and for same-support point and
bootstrap comparisons; edge-stability and paired same-support split analyses
used the primary transform.

Molecular identities for DAVIS and PKIS2 were recomputed as Standard InChIKeys
from the source SMILES with pinned RDKit.  The complete DOCKSTRING support was
aligned to the frozen raw-identity contract described above.  Exact matching
required equality of the full Standard InChIKey.  If an identity occurred more
than once on either side, each ligand--target cell was collapsed by the median
within that side before the sorted identity intersection was aligned.  Vina
scores in the matched block were clipped above at zero.  This yielded 59 DAVIS
exact identities, of which three had a constant experimental profile at
pKd=5, giving the primary 56-ligand informative block; all 154 PKIS2 exact
identities were informative.

The broad docking reference retained complete DOCKSTRING rows after excluding
every row whose raw 14-character InChIKey connectivity block occurred anywhere
in either full DAVIS or full PKIS2 panel.  This removed 254 rows and retained
259,806.  Its positive scores were clipped at zero.  The same 21 targets and
two-way-centred map definition were then used for the broad-versus-experimental
comparison.  The exact-support and broad-reference comparisons therefore
differ only in ligand-support contract, not target set or transformation.

Chemical clusters for all experimental split and bootstrap analyses were
generated from binary RDKit Morgan fingerprints using
\texttt{GetMorganGenerator(radius=2, fpSize=2048)} and
\texttt{GetFingerprint} \cite{morgan1965,rogers2010,rdkit2025}.  Pairwise
distance was one minus Tanimoto similarity.  Butina clustering used a distance
cutoff of 0.35 (Tanimoto similarity at least 0.65),
\texttt{isDistData=True} and \texttt{reordering=True} \cite{butina1999}.  For
completeness, the public loader's Murcko labels used canonical cyclic scaffolds
and grouped an acyclic molecule by its Standard-InChI connectivity block
(\texttt{ACYCLIC:<connectivity>}); those Murcko labels were not the grouping
unit in the reported experimental splits or bootstraps.  This differs from the
source-row-singleton convention used in the DOCKSTRING domain control.

Disjoint-half repeatability used 2,000 repetitions.  A ligand split randomly
permuted rows and assigned $\lfloor N/2\rfloor$ rows to each half, leaving one
row unused when $N$ was odd.  A cluster split randomly ordered complete Butina
clusters and assigned each successive cluster to the currently smaller half;
no cluster crossed halves.  Each half was mapped independently.  One complete
target-label relabeling per split supplied local sign and top-decile null
comparators.  Edge perturbation used 2,000 ordinary ligand bootstraps and 2,000
Butina-cluster bootstraps.  A cluster bootstrap sampled the number of distinct
clusters with replacement and included every member of each sampled cluster,
so its ligand count could vary.

Same-support uncertainty used paired $n$-out-of-$n$ resampling: the same ligand
or Butina-cluster draw was applied to the matched docking and experimental
matrices.  The ordinary empirical-correlation result was primary; Oracle
Approximating Shrinkage covariance converted to correlation was a finite-sample
estimator sensitivity \cite{chen2010}.  A separate paired split-half analysis
compared within-modality and cross-half maps directly.  No Spearman--Brown
extrapolation, attenuation correction or classical noise ceiling was applied,
because the nonlinear map statistic and the two assay modalities cannot be
assumed to measure one interchangeable latent quantity.

The experimental DAVIS--PKIS2 comparison used the informative 67- and
644-profile maps.  For each of 9,999 QAP permutations, one complete permutation
of the PKIS2 target labels was applied jointly to its rows and columns while the
DAVIS map remained fixed \cite{hubert1976}.  The two-sided probability was
$(1+\#\{|\rho_b|\geq|\rho_{\mathrm{obs}}|\})/(9,999+1)$.  With 210 target-pair
edges, the top positive decile contained $\lceil0.10\times210\rceil=21$ edges,
selected with deterministic stable tie handling.  QAP respects the shared-node
dependence of map edges, but remains conditional on this fixed target panel.

HotSpot repeatability used 2,000 disjoint random ligand-row halves because the
publisher workbook contains names and CAS numbers but not structures. A
separate PubChem identity audit resolved 155 of 176 records; unresolved records
were retained as singletons. The resolved standardized structures defined
radius-2, 2,048-bit Morgan fingerprints and Butina clusters at Tanimoto 0.5 for
a secondary positive-multiplier sensitivity, but not for the primary split-half
summary. Pairwise cross-assay concordance compared corresponding upper-triangle
edges of the complete DAVIS, PKIS2 and HotSpot raw or row-centred maps.

All experimental resampling used base seed 20260810, 2,000 split repetitions
and 2,000 bootstraps.  DAVIS used panel seed 20260810 and PKIS2 used 21260810
(panel stride 1,000,000).  Split seeds were
\texttt{panel\_seed+10000+counter} in the serialized loop order over full then
informative support, two-way then raw transform, and ligand then Butina unit.
For the edge bootstraps, \texttt{counter} was 4 after the full-support splits
and 8 after the informative-support splits, and \texttt{unit} was 0 for ligand
and 1 for Butina resampling; the seed was
\texttt{panel\_seed+100000+10*counter+unit}.  Paired same-support bootstrap seeds
were
\[
\texttt{panel\_seed}+200000+10000\,\texttt{support}
 +100\,\texttt{transform}+\texttt{unit},
\]
where \texttt{support}=0 for the primary informative exact matches and 1 for all
exact matches, \texttt{transform}=0 for two-way and 1 for raw, and the same unit
indices.  Paired same-support split seeds were
\[
\texttt{panel\_seed}+300000+10000\,\texttt{support}+\texttt{split},
\]
with \texttt{split}=0 for ligand and 1 for Butina.  Cross-panel QAP seeds were
$20260810+5{,}000{,}000=25260810$ for two-way-centred and 25260811 for raw
maps.  Percentile bootstrap and repeated-split ranges are
sampling/perturbation ranges for nonlinear statistics on observed supports, not
target-population confidence intervals.

\subsection{Assay-value-blind fixed-target panel transfer}

Panel transfer used the same ordered 21 targets and a Vina calibration matrix
that excluded the union of all DAVIS and PKIS2 connectivity blocks and every
conservative HotSpot identity candidate. Resolved HotSpot structures contributed
their audited RDKit connectivity; each unresolved record contributed all
plausible PubChem candidate connectivities. This removed 419 rows from the
260,060-row complete DOCKSTRING support and retained 259,641 rows. Positive
Vina scores were clipped at zero, and row centring used only the fixed 21 target
columns. Identity resolution and exclusion were performed without assay values.

For $k=4,6,8,10,12$, all $\binom{21}{k}$ panels were enumerated. Separate raw
and row-centred Vina maps selected every numerically optimal panel under each of
three distances: signed $1-r$ (primary), $1-|r|$, and $1-r^2$. Each selected
panel was then evaluated on the corresponding row-centred DAVIS, PKIS2 or
HotSpot experimental map using the same coverage loss, the mean nearest-target
distance from all 21 targets to the selected set. No experimental value entered
selection. Ties were reported lexicographically, averaged over all optima, and
conservatively as the best raw-panel loss minus the worst residual-panel loss.
Positive raw-minus-residual contrast favours the residual selector.

The primary joint summary was the mean across the three assay panels of the
normalized trapezoidal area under the contrast $c(k)$ over $k$. For the ordered
grid $k_1<\cdots<k_M$, this was
\begin{equation*}
\mathrm{AUC}_{\mathrm{norm}}
=\frac{1}{k_M-k_1}\sum_{m=1}^{M-1}
\frac{(k_{m+1}-k_m)\{c(k_m)+c(k_{m+1})\}}{2},
\end{equation*}
so normalization is by the observed $k$ range only and leaves the signed
$1-r$ loss scale unchanged. One common random
permutation of all 21 target labels was applied to every assay panel, $k$,
objective and tie statistic in each of 20,000 repetitions. One-sided
probabilities used the add-one rule. A maximum-statistic adjustment over the
three distance objectives was reported separately for each tie statistic; the
conservative signed result is primary. A sequence baseline applied the same
exact panel enumeration to the fixed 21-target global sequence-identity map.
Re-enumerating a selector after relabelling its complete Vina map is exactly
equivariant to applying that relabelling to every member of the already
enumerated numerical-optimum set. The implementation uses the latter form and
then evaluates those permuted optima against each fixed experimental distance
map; selection is therefore inside the target-label null rather than held fixed
in biological label space.

For the HotSpot chemical-support sensitivity, every member of a resolved
Morgan/Butina cluster received one shared independent $\mathrm{Exp}(1)$ weight;
the 21 unresolved compounds were singleton clusters. Two thousand weighted
maps were evaluated with the already selected Vina panels. This multiplier is a
descriptive perturbation of chemical support; target-label permutation remains
the inferential unit for panel transfer.

Target influence was assessed by omitting each kinase in turn, recomputing all
raw and row-centred maps, and repeating complete panel enumeration on the
remaining 20 targets for the same $k$ grid and primary signed objective. Thus
the selectors themselves, rather than only their evaluation edges, changed in
each leave-one-target-out analysis. Cross-assay raw and residual map agreement
was recomputed on the same reduced panels. The resulting ranges are fixed-panel
influence sensitivities, not intervals for a target superpopulation.

\subsection{Docking-44 missing-data analyses}

The primary target-mean imputation was compared with target-median imputation
on the same 12,651 rows and with a complete-case analysis on the 9,297 rows
having all 44 scores.  Restricting the primary imputed matrix to these rows is
exactly identical to their observed score matrix, making the complete-case
contrast a support contrast rather than an imputation contrast.

An observed-cell additive least-squares sensitivity estimated grand, ligand and
target effects by alternating exact unweighted updates over available cells.
The convergence tolerance was $10^{-12}$ with at most 10,000 iterations; the
released fit converged in seven.  Residuals were calculated at observed cells,
and every missing residual was set to zero, its fitted additive expectation.
This completes a residual surface for spectral calculation; it does not impute
or recover the missing raw docking scores.

To separate the unusual complete-row chemistry from its smaller sample size,
1,000 simple-random supports of 9,297 rows were drawn without replacement from
the full primary processed surface (seed 20260811).  Their residual PR
distribution is a descriptive equal-$N$ comparator, not a randomization test of
missingness.  Molecular weight, rotatable-bond count and cLogP were read from
the frozen \texttt{MW, g/mol}, \texttt{rotatable\_bonds} and \texttt{logP}
columns; heavy-atom count was recomputed with RDKit from cleaned SMILES,
falling back to canonical SMILES.  These quantities were summarized for
complete and incomplete rows, with excluded-minus-complete raw shifts,
standardized mean differences, median/IQR shifts and Cliff's delta.  No missing-
at-random mechanism was assumed and no hypothesis test assigned causality to
these associations.

\subsection{Software and computational reproducibility}

The public-core environment used Python~3.12.11, NumPy~1.26.4,
pandas~2.3.2, SciPy~1.16.2, scikit-learn~1.7.2, RDKit~2025.03.6,
Matplotlib~3.10.6 and openpyxl~3.1.5; pytest~9.1.1 ran the focused tests
and xlrd~2.0.2 read the publisher XLS
\cite{numpy2020,scipy2020,rdkit2025,matplotlib2007}.  The exact environment is
pinned in \path{requirements.lock} and \path{Dockerfile}.

The source-level public producers are:
\begin{itemize}
\item \path{analysis/cross_panel_overlap_audit.py};
\item \path{analysis/public_residual_null_audit.py};
\item \path{analysis/public_docking44_missing_data_sensitivity.py};
\item \path{analysis/public_chemical_domain_controls.py};
\item \path{analysis/public_descriptor_domain_specificity.py};
\item \path{analysis/revision_map_decision_sensitivities.py};
\item \path{analysis/public_scaffold_holdout_recovery.py};
\item \path{analysis/public_panel_selector_controls.py};
\item \path{analysis/public_panel_selector_robustness.py};
\item \path{analysis/experimental_map_reliability.py};
\item \path{analysis/fetch_anastassiadis_supplement.py};
\item \path{analysis/public_anastassiadis_identity_audit.py}; and
\item \path{analysis/public_anastassiadis_panel_validation.py}.
\end{itemize}
Running each with its recorded defaults recreates its named result directory
from the redistributed row-level sources, the checksum-fetched HotSpot workbook
and the checksum-listed public identity and target annotations. Corresponding
\texttt{test\_*.py} files lock dimensions, preprocessing, support fingerprints,
singleton handling, identity matching, seed contracts and numerical invariants.

Finally, \path{analysis/build_public_core_evidence.py} reopens the allowlisted
public sources and producer outputs, recomputes the primary spectra directly,
checks cross-artifact consistency, and writes the schema-4.0.0 ledger
\path{results/public_core_evidence.json}.  It also writes
\path{results/public_core_source_manifest.csv} and its SHA-256 digest; every
opened input and output is recorded with path, role, byte size, license and
SHA-256.  The independent \path{analysis/target_map_audit.py} command-line tool
fails closed on missing or nonnumeric target cells unless target-mean or
target-median imputation is explicitly requested, requires an absent or empty
output directory, and emits the preprocessing contract, raw and two-way maps,
spectra, signed edges and a per-run checksum manifest.  Its optional domain,
row-disjoint pilot and deterministic $1-|r|$ medoid analyses implement the
study's reusable audit; they do not validate biological mechanism, affinity
accuracy or ligand-wise target retrieval.


\section{Results}

\subsection{A shared ligand axis masks a dependent residual target map}

The two full score matrices showed the same qualitative structure despite their
different ligand collections and target panels. In Docking-44, the first
principal component explained 73.3\% of column-standardized variance and the PR
dimension was \DffRawPR{}. In DOCKSTRING-58, the corresponding values were
65.9\% and \DsRawPR{}. The PC1 loading vector was nearly parallel to a uniform
target vector (cosine similarities \DffAxisCosine{} and \DsAxisCosine{};
Fig.~\ref{fig:spectra}a), consistent with a shared favourable-score direction
rather than one target-specific contrast.

For a $P$-target correlation matrix $C$, PR dimension is
\begin{equation}
r_{\mathrm{PR}}
=\frac{\left(\sum_k\lambda_k\right)^2}{\sum_k\lambda_k^2}
=\frac{P}{1+(P-1)\overline{r_{ij}^{2}}},
\label{eq:pr}
\end{equation}
where $\lambda_k$ are the eigenvalues and the bar averages squared correlations
over unique target pairs. Thus the headline dimensions correspond directly to
mean squared target correlations of \DffRawMeanRtwo{} and \DsRawMeanRtwo{};
PR provides an interpretable common scale but no information beyond the complete
correlation spectrum.

We removed additive ligand and target main effects,
\begin{equation}
\Gamma_{ij}=X_{ij}-\bar X_{i\cdot}-\bar X_{\cdot j}+\bar X_{\cdot\cdot},
\label{eq:center}
\end{equation}
and recomputed unit-variance target correlations. Thus the raw and residual
spectra use the same post-transformation scaling. In Eq.~\ref{eq:center}, correlation already removes column
means, the practical difference from the original map is removal of the
ligand-wise common score. PR dimensions increased to \DffResidualPR{} and
\DsResidualPR{}, while mean squared correlations fell to
\DffResidualMeanRtwo{} and \DsResidualMeanRtwo{} (Fig.~\ref{fig:spectra}b,c).
The residual maps contained both positive and negative edges, rather than the
mostly positive raw maps, but they were not close to independent noise.

Row centring imposes a zero row sum and therefore a compositional constraint,
including one zero spectral mode. We tested whether that constraint or
ligand-specific residual magnitude could explain the observed concentration on
the same deterministic 12,000-ligand supports. Median residual PR was
\DffNullMedianLow{}--\DffNullMedianHigh{} in Docking-44 and
\DsNullMedianLow{}--\DsNullMedianHigh{} in DOCKSTRING under (i) an additive
Gaussian surface with target-specific residual variances, (ii) independent
permutation of empirical residual target columns, and (iii) random directions
in the zero-sum target subspace preserving every ligand's residual row norm
(Fig.~\ref{fig:spectra}e). Every null value exceeded the corresponding observed
PR. The observed residual dependence is therefore not explained by these three
specified transformation-matched null constructions. This model-conditional
comparison does not exclude every mechanical alternative, distinguish among
structured alternatives or identify the dependence as biologically correct.

A fourth null asked how much of the residual map follows from target-specific
molecular-size gradients alone. It independently permuted each processed raw
target column within ten molecular-weight bins, thereby preserving each
target's empirical score distribution within a coarse size domain while
destroying cross-target ligand correspondence, and then repeated row centring.
Median null PR fell to \DffMWConditionedNullPR{} and
\DsMWConditionedNullPR{}, explaining \DffMWConditionedExplainedShare{} and
\DsMWConditionedExplainedShare{} of the one-bin-null to observed-PR gap. Yet
the same null recovered much of the residual edge ordering (median map
Spearman \DffMWConditionedMapRho{} and \DsMWConditionedMapRho{}) and edge signs
(\DffMWConditionedSignAgreement{} and \DsMWConditionedSignAgreement{}).
Size-conditioned target marginals can therefore reproduce much of \emph{which}
target pairs co-vary, but not the strength and spectral concentration of their
joint dependence. Random partitions with the same bin sizes produced no such
shift, and five- to twenty-bin sensitivities approached a plateau
(Supplementary Section~S2.3).

Two sensitivities bounded preprocessing choices. Without target-wise variance
standardization, covariance-spectrum PR still increased from \DffCovRawPR{} to
\DffCovResidualPR{} and from \DsCovRawPR{} to \DsCovResidualPR{}; scaling
therefore affects the magnitude but not the direction of the result. In
Docking-44, full-support target-median imputation gave residual PR
\DffMedianImputeResidualPR{} and observed-cell additive fitting gave
\DffMissingAwareResidualPR{}, close to the primary \DffResidualPR{}. Complete
rows alone gave \DffCompleteCaseResidualPR{}, but the same-sized random supports
ranged only \DffEqualNResidualPRLow{}--\DffEqualNResidualPRHigh{}; the complete-
case shift was attributable to a strongly size-enriched chemical support rather
than to sample size. Restricting the mean-imputed matrix to those same complete
rows reproduced the complete-case value exactly, excluding the imputed cells as
the source of that contrast (Supplementary Section~S2).

\begin{figure}[htbp]
\centering
\includegraphics[width=\linewidth]{figures/public_core/fig1_spectra_residual.pdf}
\caption{Shared and residual spectral structure of the two public Vina panels.
(a) Target loadings on the leading score-surface principal component, oriented
toward the positive uniform vector; dotted lines mark uniform loadings and the
inset reports cosine similarity to that vector. (b) Correlation-spectrum
variance fractions before and after two-way centering; the algebraically zero
residual mode is omitted from the logarithmic display. (c) Hierarchically
ordered residual target--target correlation map for Docking-44. (d) The
corresponding residual map for DOCKSTRING-58. (e) Observed residual
participation-ratio dimension, normalized by its $P-1$ maximum, against
additive-Gaussian, column-permutation, row-norm-preserving, and ten-bin
molecular-weight-conditioned raw-score nulls on the same deterministic
$n=12{,}000$ ligand support. Null points are medians and bars are 2.5--97.5\%
simulation ranges.}
\label{fig:spectra}
\end{figure}

\subsection{Target maps vary across multiple ligand-property domains, most strongly along molecular size}

We next asked whether a residual map estimated on one part of chemical space
described another. Molecular weight was selected after preliminary inspection
and is treated as an exploratory domain coordinate, not a causal mechanism. In
the lower and upper molecular-weight quartiles, residual edge-rank agreement was
\DffMWQuartileRho{} for Docking-44 and \DsMWQuartileRho{} for
DOCKSTRING-58 (Fig.~\ref{fig:domain}a). The Docking-44 value uses all 12,651
ligands; the DOCKSTRING value uses a fixed, score-blind 15,000-ligand random support
(seed 71) for descriptor computation. This was not the disagreement expected
from finite samples of the same library: matched-size random-disjoint supports
gave agreements near one. Nor was it a generic consequence of chemical
disjointness. When whole Butina clusters (Docking-44) or Bemis--Murcko groups
(DOCKSTRING-58) were assigned to disjoint halves while matching the molecular-
weight distribution, median residual-map agreements were
\DffMWMatchedControlRho{} and \DsMWMatchedControlRho{}.

The support effect was graded rather than unique to one molecular axis. Using
inclusive lower/upper threshold tails, molecular weight, heavy-atom
count and Labute accessible surface area gave residual-map agreements of
\DffSizeAxisRhoRange{} in Docking-44 and \DsSizeAxisRhoRange{} in
DOCKSTRING. The corresponding ranges across topological polar surface area,
clogP, rotatable-bond count and ring count were \DffOtherAxisRhoRange{} and
\DsOtherAxisRhoRange{} (Fig.~\ref{fig:domain}d). Each of the 14 observed
contrasts lay below the 2.5th percentile of 200 random-row-disjoint controls
matched to its actual low/high tail sizes. Molecular weight correlated
strongly with heavy-atom count ($\rho=0.935$ and 0.965) and accessible surface
area ($\rho=0.968$ and 0.979), identifying one highly collinear size-related
family rather than three independent replications. All seven comparisons were
examined post hoc. The gradient in effect size narrows the observed association
but does not establish size or contact opportunity as causal; the four other
axes also differed from their matched random controls. Inclusive thresholds
retained ties, so discrete descriptors produced larger tails (the DOCKSTRING
ring-count split covered the entire 15,000-row support between its two disjoint
tails).

The contrast was not tied to the quartile cut. Across lower/upper tail fractions
from 10\% through 40\%, residual-map agreement ranged from 0.201 to 0.259 in
Docking-44 and from 0.174 to 0.528 in DOCKSTRING-58. At every threshold the
observed agreement was below the 2.5th percentile of its equal-$N$ random-
disjoint control (Fig.~\ref{fig:domain}b). Within either the low- or high-weight
band, chemical-group-disjoint halves remained highly concordant, making unstable
estimation inside a restricted band an unlikely explanation.

Finally, ten non-overlapping, equally populated molecular-weight bins replaced a
binary split. Map dissimilarity increased steadily with separation of bin
medians: the descriptive slopes were 0.220 and 0.316 dissimilarity units per
100~Da, and the corresponding Spearman correlations were
\DffMWContinuumRho{} and \DsMWContinuumRho{} (Fig.~\ref{fig:domain}c). These
45 pairwise distances reuse ten maps and are not independent observations. In
DOCKSTRING, the quartile result also reproduced across six separately seeded
15,000-ligand supports (residual $\rho$ range 0.325--0.361;
Supplementary Section~S3). Together, these controls document a large and
threshold-robust domain shift within both source matrices; they do not establish
molecular weight as its sole cause.

\begin{figure}[htbp]
\centering
\includegraphics[width=\linewidth]{figures/public_core/fig2_support_domain.pdf}
\caption{Chemical-domain dependence of residual Vina target-correlation maps.
(a) Diamonds show the observed edge-rank Spearman correlation between low- and
high-molecular-weight quartile maps (0.205 for Docking-44 and 0.351 for
DOCKSTRING-58). Open circles and bars show medians and 2.5--97.5\% ranges across
200 whole-chemical-group-disjoint splits for the full support after MW matching
and separately within each MW quartile; the global matched medians were 0.992
and 0.994. Chemical groups are inherited source cluster labels for Docking-44 and
Bemis--Murcko scaffolds, with acyclic compounds as singletons, for
DOCKSTRING-58. The small DOCKSTRING ticks show the observed residual contrast
in six seeded 15,000-ligand supports; Docking-44 uses all 12,651 ligands. These
repeated-split composition ranges and seeded-support values are not confidence
intervals. (b) Residual-map agreement across 10--40\% low/high tail definitions;
dotted lines and bands are means and 2.5--97.5\% ranges across 100 equal-size
random-row-disjoint control supports; unlike panel a, chemical groups may recur
across their halves. The quartile definition differs from the 25\% sweep point
by at most one boundary ligand because the producers use different tie rules.
(c) Pairwise map dissimilarity among ten equal-size MW bins versus their
median-MW separation. Marker size encodes bin-index separation; the fits are
descriptive because the 45 pairs reuse ten maps. (d) Residual low/high
threshold-tail map agreement along seven ligand-descriptor axes. Filled points
are observed contrasts; open points and bars are medians and 2.5--97.5\%
ranges from 200 matched-tail-size random-row-disjoint controls. Every observed
contrast lies below its control range, with the largest changes along the
correlated molecular-size axes. Inclusive thresholds retain ties, so discrete
descriptors can have expanded tails. All descriptor comparisons were selected
post hoc and are descriptive; these panels diagnose support dependence but do
not identify a causal molecular property or establish biological target
relationships.}
\label{fig:domain}
\end{figure}

\subsection{Source-library pilots recover target maps, but residual scores resist panel compression}

Support dependence implies that a map should be calibrated on the intended
library, but not necessarily that the complete screen must be run first. We drew
100 simple-random pilot samples of 200 or 500 ligands from each source matrix and
compared each estimated residual map with the map on its row-disjoint complement.
With 200 ligands, mean edge-rank agreement was \DffProbeTwoHundred{} for
Docking-44 and \DsProbeTwoHundred{} for DOCKSTRING-58; with 500 it was
\DffProbeFiveHundred{} and \DsProbeFiveHundred{}
(Fig.~\ref{fig:recovery}a). Edge-sign and strongest-edge recovery were lower
than global agreement, and are reported in Supplementary Section~S4.

Recovery also persisted under five-fold chemical-group holdout. Calibration
ligands were sampled from four folds and compared with the residual map of the
fifth; cyclic scaffolds had zero overlap by construction and acyclic compounds
were separated by standardized connectivity. Mean agreements at 200 ligands
were \DffScaffoldProbeTwoHundred{} and \DsScaffoldProbeTwoHundred{}, increasing
to \DffScaffoldProbeFiveHundred{} and \DsScaffoldProbeFiveHundred{} at 500.
Matched random-row holdouts were similar, and all paired central ranges for the
group-minus-random contrast included zero. This is a same-library
chemical-group sensitivity, not transport to a shifted external library.

Map recovery alone did not establish that a small panel preserved the score
surface. We therefore used five chemical-group-disjoint folds, selected
$k=4,6,8,10,12$ targets from a 500-ligand calibration sample, fitted
multivariate Ridge models using only those selected raw target scores, and
predicted omitted targets in the held-out fold. For the raw score surface,
median variance-weighted omitted-target $R^2$ increased rapidly with $k$; at
$k=8$ it was \DffSelectorRawRtwo{} in Docking-44 and \DsSelectorRawRtwo{} in
DOCKSTRING-58. For the row-centred residual surface, however, the corresponding
values were only \DffSelectorResidualRtwo{} and
\DsSelectorResidualRtwo{} (Fig.~\ref{fig:recovery}b). Thus most apparent raw
compressibility followed the shared ligand axis removed by centring, whereas
target-relative deviations were distributed across more targets.

This residual gap was not a bookkeeping artifact caused by predicting raw
scores and centring them afterwards. Directly fitting the residual outcomes
from the same selected raw predictors gave $R^2=\DffSelectorDirectResidualRtwo{}$
and \DsSelectorDirectResidualRtwo{}, differing from the derived-residual route
by only \DffSelectorDirectResidualDelta{} and
\DsSelectorDirectResidualDelta{}. Both multivariate fits exceeded a
selected-target PC1 baseline (\DffSelectorResidualPCOneRtwo{} and
\DsSelectorResidualPCOneRtwo{}; Fig.~\ref{fig:recovery}d). On the targets
omitted by both a designed panel and each of three fixed random panels, the
designed residual panel improved mean $R^2$ in every fold; median gains were
\DffSelectorCommonResidualDelta{} and \DsSelectorCommonResidualDelta{}.
Because only three random panels were fixed in advance, this comparison is
descriptive rather than a tail probability.

The simpler map-coverage objective remained useful but answered a weaker
question. Eight medoids selected from a 200-ligand pilot retained 72\% and 80\%
of the full-map improvement over median random panels, rising to 85\% and 90\%
at 500 ligands (Fig.~\ref{fig:recovery}c). A two-stage design with 200 fully
scored pilot ligands and eight targets thereafter would use 19.5\% and 13.9\%
of the complete score cells. These savings describe Vina computation, not
affinity prediction, and the residual reconstruction results show why raw-map
coverage must not be mistaken for preservation of target-relative information.

\begin{figure}[htbp]
\centering
\includegraphics[width=\linewidth]{figures/public_core/fig3_pilot_decision.pdf}
\caption{Recovery and compression of source-library Vina maps. (a) Global
edge-rank agreement between maps estimated from 200 or 500 random calibration
ligands and their row-disjoint complement maps. Points are means and bars are
2.5--97.5\% ranges over 100 supports. (b) Median omitted-target
variance-weighted $R^2$ for score and row-centred residual surfaces when
$k=4$--12 selected raw target scores are observed. Lines join medians and bars
span five chemical-group-disjoint folds. (c) Pilot-medoid coverage expressed as
the fraction of the random-to-full-map nearest-distance gap remaining. These
are Vina score-compression diagnostics, not affinity predictions or a biological
target optimum. (d) At $k=8$, multivariate Ridge residual reconstruction after
raw-score prediction, direct residual-outcome Ridge using the same selected raw
predictors, and a selected-target PC1 baseline. Points are fold medians and bars
span five folds; the direct fit is an artifact check rather than a deployment
policy.}
\label{fig:recovery}
\end{figure}

\subsection{Residual-map selection improves fixed-target co-response coverage without validating affinity prediction}

We first established that the experimental target maps were themselves
repeatable. On the same fixed 21 kinases, median disjoint-half residual-map
agreement was \DavisSplitRho{} for DAVIS, \PkistwoSplitRho{} for PKIS2 and
\HotspotSplitRho{} for the \HotspotLigands{}-compound Anastassiadis HotSpot
panel (Fig.~\ref{fig:experimental}a). The HotSpot 2.5--97.5\% split range was
\HotspotSplitRhoLow{}--\HotspotSplitRhoHigh{}. These panels differ in compounds
and assay format and are not independent samples of a common population, but
they supply three observed co-response geometries on the same target labels.

Removing each compound's mean response made those geometries more concordant.
DAVIS--PKIS2 edge-rank agreement increased from
\ExperimentalCrossRawRho{} to \ExperimentalCrossRho{}; DAVIS--HotSpot increased
from \DavisHotspotRawRho{} to \DavisHotspotResidualRho{}, and PKIS2--HotSpot
from \PkistwoHotspotRawRho{} to \PkistwoHotspotResidualRho{}
(Fig.~\ref{fig:experimental}d). This does not show that centring recovers a true
biophysical network. It shows that a target-relative co-response estimand is
more reproducible across these three assay contexts than the corresponding raw
correlation map. The residual-over-raw ordering survived all 21
leave-one-target-out omissions for every assay pair. Because DAVIS contains a
substantial released pKd floor, these cross-assay arrows remain descriptive;
the panel-transfer result below was also checked after independently ranking
ligands within each experimental target before row centring. That scale-free
sensitivity retained \PanelRankPositiveCells{}/15 positive cells (mean
contrast \PanelRankConservativeMean{}, adjusted $p=\PanelRankP{}$).

We then tested a bounded design consequence without using assay values for
selection. For each of $k=4,6,8,10,12$, all target subsets were scored on either
the raw or row-centred DOCKSTRING map after every plausible DAVIS, PKIS2 or
HotSpot evaluation-compound connectivity had been removed from the calibration
support. Numerically optimal subsets were evaluated by signed experimental
co-response coverage loss $1-r$ on each assay panel. To avoid favourable tie
selection, the primary contrast used the best raw-map optimum minus the worst
residual-map optimum; positive values favour the residual selector. This
conservative contrast was positive in \PanelTransferPositiveCells{} of
\PanelTransferTotalCells{} assay--$k$ cells, with mean
\PanelTransferConservativeMean{} (Fig.~\ref{fig:experimental}b). The joint
normalized area-under-the-loss-curve improvement was \PanelTransferJointAUC{};
a shared 20,000-permutation target-label test, adjusted for the maximum over the
three assay objectives, gave one-sided $p=\PanelTransferJointP{}$. Because the
three-panel transfer analysis was introduced post hoc, this is an exploratory
conditional randomization result on the fixed 21-target set, not confirmatory
error control for the full analysis history. In a HotSpot
chemical-cluster multiplier, all 2,000 perturbations favoured the residual
selector (median \HotspotClusterMedian{}, 2.5--97.5\%
\HotspotClusterLow{}--\HotspotClusterHigh{}).
The result was not carried by one kinase: after omitting each target and
reselecting both panels on the remaining 20-target maps, the joint AUC contrast
remained positive in \PanelTargetLooPositive{}/\PanelTargetLooTotal{} analyses
(range \PanelTargetLooAucLow{}--\PanelTargetLooAucHigh{}), with
\PanelTargetLooCellWinsLow{}--\PanelTargetLooCellWinsHigh{} of 15 cells
positive per omission. These strongly overlapping leave-one-target-out values
are influence checks, not 21 independent confirmations or a second $p$-value;
their magnitude is small relative to the full $0$--$2$ signed-distance scale.

Sequence identity remained a competitive panel-design comparator: the
lexicographic residual Vina panel had lower loss in only 9 of the 15 cells.
Accordingly, the result is not evidence that docking replaces sequence-based
panel design. It shows that, when Vina is used as the design signal, removing
the shared ligand axis changes the selected subset in a direction that more
often covers observed target-relative co-response structure on this fixed
kinase panel.

The affinity boundary remained deliberately separate. Exact full Standard
InChIKey matching to DOCKSTRING yielded \DavisSameSupportLigands{} informative
DAVIS ligands and \PkistwoSameSupportLigands{} PKIS2 ligands. Vina--experiment
residual-map agreement was only \DavisSameSupportRho{} and
\PkistwoSameSupportRho{} (Fig.~\ref{fig:experimental}c), with broad de-leaked
DOCKSTRING references of \DavisBroadSupportRho{} and
\PkistwoBroadSupportRho{}. Chemical-cluster bootstrap ranges were 0.098--0.369
and 0.163--0.391. Thus panel-level transfer does not imply that Vina predicts
affinities, compound orderings or experimental maps on identical ligands.

Nor is a target map a ligand-wise ranking rule. Subtracting a ligand row mean is
exactly order invariant:
\begin{equation}
(X_{ij}-\bar X_{i\cdot})-(X_{ik}-\bar X_{i\cdot})=X_{ij}-X_{ik}.
\label{eq:ranking}
\end{equation}
Equation~\ref{eq:ranking} shows that only a target-dependent operation, such as unequal target scaling, can alter
that order, in which case it defines a new score rule requiring task-level
validation. This algebraic boundary prevents the increase in residual dimension
from being read as evidence of improved target fishing.

\begin{figure}[htbp]
\centering
\includegraphics[width=\linewidth]{figures/public_core/fig4_experimental_boundary.pdf}
\caption{Experimental-map reproducibility, assay-value-blind panel transfer,
and the docking--experiment boundary on 21 fixed kinase targets. (a)
Two-way-centred experimental edge-map agreement across 2,000 disjoint halves of
DAVIS, PKIS2 and the Anastassiadis HotSpot panel. DAVIS and PKIS2 keep Butina
clusters within a half; HotSpot uses ligand rows. (b) Signed co-response
coverage-loss contrasts over $k$. Solid circles are the conservative best-raw
minus worst-residual contrast over every numerically optimal Vina panel;
positive values favour the residual selector. Dotted crosses compare the
sequence-identity panel with the lexicographic residual panel. Vina panels were
selected without assay values after excluding all plausible evaluation-compound
connectivities from the calibration support. (c) Score- and residual-map
agreement between Vina and experiment on 56 exact DAVIS and 154 exact PKIS2
Standard InChIKey matches, compared with de-leaked broad-support Vina maps;
bars are 2.5--97.5\% ranges across 2,000 Butina-cluster bootstraps. (d) Paired
raw and row-centred edge-map agreement across DAVIS--PKIS2, DAVIS--HotSpot and
PKIS2--HotSpot assay pairs. All comparisons are conditional on the fixed target
panel; cross-panel experimental agreement is not an attenuation-corrected noise
ceiling.}
\label{fig:experimental}
\end{figure}

\section{Discussion}

The main result is a distinction between two kinds of apparent redundancy. Raw
Vina matrices were dominated by a shared ligand axis and were readily
reconstructed from a small target subset. After that axis was removed, the
remaining target-relative surface had higher effective dimension and markedly
lower held-out reconstruction accuracy. The increase in PR therefore has a
concrete meaning: a reduced panel that preserves general score favourability can
still discard much of the variation that distinguishes targets from one
another. It does not mean that the residual is biologically correct, that
centring improves docking, or that every target must be measured.

The novelty is consequently not PR, PCA or two-way centring. PR is an exact
re-expression of mean squared target correlation (Eq.~\ref{eq:pr}), and docking
profile normalization predates this study
\cite{fukunishi2005,fukunishi2006,sepresa2009,luo2017}. The contribution is a
transformation-matched audit that connects spectrum, ligand support and a
declared panel-design loss. It shows why map recovery, score reconstruction and
experimental panel coverage must be tested separately rather than inferred from
one low-dimensional summary.

That separation changes the interpretation of computational savings. A few
hundred source-library ligands recovered global map ordering and a medoid
coverage objective, but the stronger out-of-group test showed that eight targets
explained 70--90\% of raw omitted-target variance and only 22--28\% of residual
variance. Direct residual fitting produced the same result, excluding a
prediction-then-centring artifact. Designed residual panels also outperformed
three fixed random panels on common omitted targets in every fold. These are
score-surface results on the tested ligand domains; they are not affinity
prediction or permission to omit mandatory safety targets.

The three experimental panels provide a different, bounded consequence.
Row-centred co-response maps were more concordant than raw maps for every assay
pair. Without using assay values, residual-map Vina panels had lower signed
co-response coverage loss than raw-map panels in 14 of 15 assay--size cells,
including a shared-label permutation result across DAVIS, PKIS2 and HotSpot;
reselection after each single-target omission preserved a positive joint
contrast in all 21 cases.
This supports residual Vina maps as one panel-design signal on the fixed 21
kinases. Sequence identity remained competitive, and exact-support
Vina--experiment correlations were only moderate. The analysis therefore does
not establish affinity calibration, target retrieval, compound ranking or the
superiority of docking over sequence-based design.

Ligand support remains part of every map claim. Residual maps changed across
all seven examined descriptor domains, most strongly along the collinear
molecular-weight, heavy-atom and accessible-surface family. Three matched nulls,
including preservation of each residual row norm, showed that the residual
concentration was not a consequence of the centring constraint alone. A
size-conditioned null recovered much of the signed edge ordering but only part
of the concentration gap. These post-hoc results identify support-dependent,
non-independent score variation, not a biochemical mechanism.

\begin{table}[htbp]
\centering
\caption{A decision-oriented audit for a dense ligand--target score matrix.}
\label{tab:audit}
\small
\begin{tabular}{p{0.06\linewidth}p{0.31\linewidth}p{0.52\linewidth}}
\toprule
Step & Operation & Decision or boundary \\
\midrule
1 & Declare the score, receptor and ligand-support contract & Prevent maps from
different populations or preprocessing rules being treated as the same object. \\
2 & Inspect the raw spectrum and PC1 loadings & Quantify a shared score axis;
do not interpret PR as algebraic rank or docking quality. \\
3 & Remove additive main effects and compare with matched nulls & Test residual
dependence beyond centring mechanics, without assigning biological meaning. \\
4 & Re-estimate the map across relevant chemical domains & Detect a support
change before transporting edges or clusters to an intended library. \\
5 & Measure row-disjoint pilot recovery at candidate sample sizes & Choose the
calibration size for the actual output of interest: global map, edge signs or a
declared panel loss. \\
6 & Reconstruct both raw and residual held-out score surfaces & Detect panels
that preserve the shared ligand axis but lose target-relative variation. \\
7 & Evaluate experimental panel utility without assay-value leakage & Keep
computational coverage, experimental co-response and affinity prediction as
separate claims. \\
\bottomrule
\end{tabular}
\end{table}

Several limits determine how far these results can travel. Both large matrices
use Vina-family scores and overlap by \SharedExactLigands{} exact molecular
identities, \SharedConnectivityBlocks{} connectivity blocks,
\SharedTargetIdentities{} mapped targets and \SharedReceptorStructures{}
receptor structures; they are separately generated resources, not independent
samples of all chemistry and proteins. Targets were fixed and all ranges are
conditional on them. Pilot calibration remained within each source library.
The three kinase panels share target identities and some compounds, use
different assay formats, and are not independent biological replications.
HotSpot inference used row rather than chemical-cluster split halves, although
the panel-transfer result retained its sign in the reported cluster multiplier.
Single-target deletion cannot exclude joint influence by a redundant kinase
group.
The sequence comparator used global identity and was not an exhaustive
protein-panel design benchmark.

Finally, the analysis begins with frozen published scores. It reruns the
matrix-level transformations and comparisons but not receptor preparation,
search-box selection, pose generation or upstream docking. It can therefore
diagnose dependence, support sensitivity and compression in the released
scores, but cannot adjudicate pose correctness, absolute affinity or docking
quality.

\section{Conclusions}

In two large Vina-family matrices, a dominant ligand axis made raw scores appear
far more compressible than the target-relative residual surface. Residual maps
were support dependent, recoverable from modest source-library pilots and useful
for choosing fixed kinase panels that more often covered signed experimental
co-response structure than raw-map panels. Exact-support affinity agreement
remained moderate and sequence-based panels remained competitive. The practical
workflow is therefore bounded: report raw and residual spectra, test
reconstruction on held-out chemical groups, select panels under a declared loss,
and validate experimental uses separately. This distinguishes computational
score compression from biological or affinity claims.

\section*{Abbreviations}

ARI, adjusted Rand index; AUC, area under the curve; GCV, generalized
cross-validation; MIP, mixed-integer programming; MW, molecular weight; OAS,
Oracle Approximating Shrinkage; PAM, Partitioning Around Medoids; PC1, first
principal component; PR, participation ratio; QAP, quadratic assignment
procedure; $R^2$, coefficient of determination; RMSE, root-mean-square error;
SSE, sum of squared errors; SST, total sum of squares.

\section*{Declarations}

\subsection*{Ethics approval and consent to participate}

Not applicable.  This secondary computational study used public molecular,
docking-score and biochemical-assay datasets and involved no human participants,
identifiable human data, human tissue or new animal experiments.

\subsection*{Consent for publication}

Not applicable.

\subsection*{Availability of data and materials}

All row-level inputs required for the submitted public-core claims are publicly
available and checksum identified; they are redistributed in the analysis
release where source terms permit.
Docking-44 is the frozen project export associated with Nikitin et al.
\cite{nikitin2025} and is included under CC~BY~4.0.  DOCKSTRING v1 and the
derived row-aligned identity contract retain Apache-2.0 and cite the source
release \cite{dockstring2022}.  The DAVIS-Complete V3 table is the exact upstream
TSV export, recompressed without content changes, and is CC0-1.0
\cite{davis2011,wu2025davis}.  The official PKIS2 S4 workbook is included under
CC~BY~4.0 \cite{drewry2017}. The official Anastassiadis Supplementary Table~3
workbook is fetched from the publisher through a SHA-256-gated script and is
not redistributed; the package includes only the identity/provenance crosswalk
and derived analysis outputs \cite{anastassiadis2011}. Third-party data retain their upstream licences;
\path{data_manifest.csv} and \path{DATA_LICENSES.md} are the machine-readable
and human-readable authorities for source version, checksum and redistribution
terms.

The MIT-licensed v4 analysis code, frozen inputs, tests, result bundles, exact
manuscript and Supplement sources, generated figures and strict public-core
source manifest have been assembled as one release candidate.  A
version-specific immutable archive will be cited only after that candidate is
tagged, rebuilt from a clean clone and deposited.  The present draft therefore
does not cite the mutable repository branch or an earlier-version DOI as its
evidence package.  Author-generated summaries and target annotations will be
released under CC~BY~4.0 unless an identified upstream condition applies.  The
candidate package permits source-level reconstruction of every central
statistical claim from Docking-44, DOCKSTRING, DAVIS, PKIS2 and the publicly
fetched HotSpot workbook without registration, a private file or a historical
aggregate ledger. It does not include or claim to rerun
upstream receptor preparation, pose generation or docking; those stages remain
within the cited Docking-44 and DOCKSTRING source studies and releases.

\subsection*{Competing interests}

The authors declare no financial or commercial competing interests.  V.S. and
M.F. are coauthors of the source study from which the Docking-44 matrix was
obtained \cite{nikitin2025}; that relationship is disclosed because the present
study reanalyses the matrix.

\subsection*{Funding}

This research received no specific grant from any funding agency in the public,
commercial or not-for-profit sectors.

\subsection*{Authors' contributions}

A.L. conceived the study, performed the analyses and drafted the manuscript.
V.S. and M.F. contributed to study design and interpretation.  All authors
reviewed the manuscript and approved the final submitted version.

\subsection*{Acknowledgements}

The authors thank the developers, curators and maintainers of the public
datasets and open-source software used in this study.

\subsection*{Use of generative artificial intelligence}

No artificial-intelligence system is listed as an author or was treated as a
source of scientific evidence.  OpenAI Codex and Anthropic Claude were used for
code review, literature-search assistance and language editing.  The authors
inspected the underlying data and code, verified the reported claims and
citations, and retain full responsibility for the work.

\subsection*{Authors' information}

ORCID: Adeliya Leleytner, 0009-0001-8214-2496; Victor Safronov,
0009-0003-5668-0201; Maxim Fedorov, 0000-0003-3901-3565.

% EDITORIAL VERIFICATION FLAGS (comments, not submission prose):
% 1. Confirm the author-supplied funding, competing-interest and contribution
%    statements with all authors immediately before submission.
% 2. Confirm all three ORCIDs and the corresponding-author metadata against the
%    journal submission system.
% 3. references_public.tex deliberately has no verified version-specific v4
%    Zenodo DOI.  Add the DOI, archive version and deposit date only after the
%    final public-core archive is deposited and its files/checksums are checked;
%    do not reuse CITATION.cff metadata from an earlier analytical release.
% 4. Confirm the final supplementary-file name, section/figure counts and licence
%    before adding an ``Additional file'' declaration; none is asserted here.
% 5. Confirm the journal's required wording/category for the generative-AI
%    disclosure while preserving the explicit no-AI-authorship statement.
% 6. After the clean-clone checks, cite the immutable v4 tag/DOI and confirm that
%    its public_core_source_manifest.sha256 matches the archive.  Do not restore
%    a mutable-branch citation or reuse a DOI from an earlier analytical version.


\begin{thebibliography}{99}
\small

\bibitem{vina2010} Trott O, Olson AJ (2010) AutoDock Vina: improving the speed and
accuracy of docking with a new scoring function, efficient optimization, and
multithreading. J Comput Chem 31:455--461. doi:10.1002/jcc.21334.

\bibitem{ding2023} Ding J, Tang S, Mei Z et al (2023) Vina-GPU 2.0: further
accelerating AutoDock Vina and its derivatives with graphics processing units.
J Chem Inf Model 63:1982--1998. doi:10.1021/acs.jcim.2c01504.

\bibitem{dockstring2022} Garc\'ia-Orteg\'on M, Simm GNC, Tripp AJ,
Hern\'andez-Lobato JM, Bender A, Bacallado S (2022) DOCKSTRING: easy molecular
docking yields better benchmarks for ligand design. J Chem Inf Model
62:3486--3502. doi:10.1021/acs.jcim.1c01334.

\bibitem{nikitin2025} Nikitin I, Morgunov I, Safronov V, Kalyuzhnaya A,
Fedorov M (2025) Towards explainable computational toxicology: linking
antitargets to rodent acute toxicity. Pharmaceutics 17:1573.
doi:10.3390/pharmaceutics17121573.

\bibitem{pan2003} Pan Y, Huang N, Cho S, MacKerell AD Jr (2003) Consideration of
molecular weight during compound selection in virtual target-based database
screening. J Chem Inf Comput Sci 43:267--272. doi:10.1021/ci020055f.

\bibitem{vigers2004} Vigers GPA, Rizzi JP (2004) Multiple active site corrections
for docking and virtual screening. J Med Chem 47:80--89. doi:10.1021/jm030161o.

\bibitem{jacobsson2006} Jacobsson M, Karl\'en A (2006) Ligand bias of scoring
functions in structure-based virtual screening. J Chem Inf Model 46:1334--1343.
doi:10.1021/ci050407t.

\bibitem{carta2007} Carta G, Knox AJS, Lloyd DG (2007) Unbiasing scoring
functions: a new normalization and rescoring strategy. J Chem Inf Model
47:1564--1571. doi:10.1021/ci600471m.

\bibitem{sepresa2009} Yang L, Luo H, Chen J, Xing Q, He L (2009) SePreSA: a
server for the prediction of populations susceptible to serious adverse drug
reactions implementing the methodology of a chemical--protein interactome.
Nucleic Acids Res 37:W406--W412. doi:10.1093/nar/gkp312.

\bibitem{casey2009} Casey FP, Pihan E, Shields DC (2009) Discovery of small
molecule inhibitors of protein--protein interactions using combined ligand and
target score normalization. J Chem Inf Model 49:2708--2717.
doi:10.1021/ci900294x.

\bibitem{luo2017} Luo Q, Zhao L, Hu J, Jin H, Liu Z, Zhang L (2017) The scoring
bias in reverse docking and the score normalization strategy to improve success
rate of target fishing. PLoS One 12:e0171433.
doi:10.1371/journal.pone.0171433.

\bibitem{lapillo2019} Lapillo M, Tuccinardi T, Martinelli A, Macchia M,
Giordano A, Poli G (2019) Extensive reliability evaluation of docking-based
target-fishing strategies. Int J Mol Sci 20:1023.
doi:10.3390/ijms20051023.

\bibitem{lee2020crds} Lee A, Kim D (2020) CRDS: consensus reverse docking system
for target fishing. Bioinformatics 36:959--960.
doi:10.1093/bioinformatics/btz656.

\bibitem{krause2023} Krause F, Voigt K, Di Ventura B, {\"O}zt{\"u}rk MA (2023)
ReverseDock: a web server for blind docking of a single ligand to multiple
protein targets using AutoDock Vina. Front Mol Biosci 10:1243970.
doi:10.3389/fmolb.2023.1243970.

\bibitem{fukunishi2005} Fukunishi Y, Mikami Y, Nakamura H (2005) Similarities
among receptor pockets and among compounds: analysis and application to
in silico ligand screening. J Mol Graph Model 24:34--45.
doi:10.1016/j.jmgm.2005.04.004.

\bibitem{fukunishi2006} Fukunishi Y, Mikami Y, Takedomi K, Yamanouchi M,
Shima H, Nakamura H (2006) Classification of chemical compounds by
protein--compound docking for use in designing a focused library. J Med Chem
49:523--533. doi:10.1021/jm050480a.

\bibitem{fukunishi2018} Fukunishi Y, Yamashita Y, Mashimo T, Nakamura H (2018)
Prediction of protein--compound binding energies from known activity data:
docking-score-based method and its applications. Mol Inform 37:e1700120.
doi:10.1002/minf.201700120.

\bibitem{lee2012profiles} Lee M, Kim D (2012) Large-scale reverse docking
profiles and their applications. BMC Bioinformatics 13(Suppl 17):S6.
doi:10.1186/1471-2105-13-S17-S6.

\bibitem{martin2011pqsar} Martin E, Mukherjee P, Sullivan D, Jansen J (2011)
Profile-QSAR: a novel meta-QSAR method that combines activities across the
kinase family to accurately predict affinity, selectivity, and cellular
activity. J Chem Inf Model 51:1942--1956. doi:10.1021/ci1005004.

\bibitem{bembenek2018} Bembenek SD, Hirst G, Mirzadegan T (2018)
Determination of a focused mini kinase panel for early identification of
selective kinase inhibitors. J Chem Inf Model 58:1434--1440.
doi:10.1021/acs.jcim.8b00222.

\bibitem{zhang2019informer} Zhang H, Ericksen SS, Lee CP et al (2019)
Predicting kinase inhibitors using bioactivity matrix derived informer sets.
PLoS Comput Biol 15:e1006813. doi:10.1371/journal.pcbi.1006813.

\bibitem{laufkotter2020} Laufk\"otter O, Laufer S, Bajorath J (2020)
Identifying representative kinases for inhibitor evaluation via systematic
analysis of compound-based target relationships. Eur J Med Chem 204:112641.
doi:10.1016/j.ejmech.2020.112641.

\bibitem{mangione2019cando} Mangione W, Samudrala R (2019) Identifying protein
features responsible for improved drug repurposing accuracies using the CANDO
platform: implications for drug design. Molecules 24:167.
doi:10.3390/molecules24010167.

\bibitem{kangas2014active} Kangas JD, Naik AW, Murphy RF (2014) Efficient
discovery of responses of proteins to compounds using active learning. BMC
Bioinformatics 15:143. doi:10.1186/1471-2105-15-143.

\bibitem{davis2011} Davis MI, Hunt JP, Herrgard S et al (2011) Comprehensive
analysis of kinase inhibitor selectivity. Nat Biotechnol 29:1046--1051.
doi:10.1038/nbt.1990.

\bibitem{wu2025davis} Wu MH (2025) A Complete and Modification-Aware DAVIS
Dataset. Harvard Dataverse, V3. doi:10.7910/DVN/RTQGP1.

\bibitem{drewry2017} Drewry DH, Wells CI, Andrews DM et al (2017) Progress
towards a public chemogenomic set for protein kinases and a call for
contributions. PLoS One 12:e0181585. doi:10.1371/journal.pone.0181585.

\bibitem{anastassiadis2011} Anastassiadis T, Deacon SW, Devarajan K,
Ma H, Peterson JR (2011) Comprehensive assay of kinase catalytic activity
reveals features of kinase inhibitor selectivity. Nat Biotechnol 29:1039--1045.
doi:10.1038/nbt.2017.

\bibitem{mazzucato2016} Mazzucato L, Fontanini A, La Camera G (2016) Stimuli
reduce the dimensionality of cortical activity. Front Syst Neurosci 10:11.
doi:10.3389/fnsys.2016.00011.

\bibitem{roy2007} Roy O, Vetterli M (2007) The effective rank: a measure of
effective dimensionality. In: 15th European Signal Processing Conference,
Pozna\'n, pp 606--610. doi:10.5281/zenodo.40328.

\bibitem{hubert1976} Hubert L, Schultz J (1976) Quadratic assignment as a
general data analysis strategy. Br J Math Stat Psychol 29:190--241.
doi:10.1111/j.2044-8317.1976.tb00714.x.

\bibitem{chen2010} Chen Y, Wiesel A, Eldar YC, Hero AO (2010) Shrinkage
algorithms for MMSE covariance estimation. IEEE Trans Signal Process
58:5016--5029. doi:10.1109/TSP.2010.2053029.

\bibitem{murcko1996} Bemis GW, Murcko MA (1996) The properties of known drugs.
1. Molecular frameworks. J Med Chem 39:2887--2893.
doi:10.1021/jm9602928.

\bibitem{morgan1965} Morgan HL (1965) The generation of a unique machine
description for chemical structures---a technique developed at Chemical
Abstracts Service. J Chem Doc 5:107--113. doi:10.1021/c160017a018.

\bibitem{rogers2010} Rogers D, Hahn M (2010) Extended-connectivity fingerprints.
J Chem Inf Model 50:742--754. doi:10.1021/ci100050t.

\bibitem{butina1999} Butina D (1999) Unsupervised data base clustering based on
Daylight's fingerprint and Tanimoto similarity: a fast and automated way to
cluster small and large data sets. J Chem Inf Comput Sci 39:747--750.
doi:10.1021/ci9803381.

\bibitem{rdkit2025} Landrum G et al (2025) RDKit: open-source cheminformatics
software, release 2025.03.6. Zenodo. doi:10.5281/zenodo.16996017.

\bibitem{numpy2020} Harris CR, Millman KJ, van der Walt SJ et al (2020) Array
programming with NumPy. Nature 585:357--362.
doi:10.1038/s41586-020-2649-2.

\bibitem{scipy2020} Virtanen P, Gommers R, Oliphant TE et al (2020) SciPy 1.0:
fundamental algorithms for scientific computing in Python. Nat Methods
17:261--272. doi:10.1038/s41592-019-0686-2.

\bibitem{matplotlib2007} Hunter JD (2007) Matplotlib: a 2D graphics environment.
Comput Sci Eng 9:90--95. doi:10.1109/MCSE.2007.55.

\end{thebibliography}


\end{document}
